



\medskip


\centerline{\large\bf Lemma on Generic Flatness}

\medskip


\begin{lemma} \label{nhq-algebraic-generic-flatness}\dcref{4.1}\source{nitsure-hilbert-quot:page-0018}
Let $A$ be a noetherian domain,
and $B$ a finite type $A$ algebra. Let $M$ be a finite $B$-module.
Then there exists an $f\in A$, $f\ne 0$, such that
the localisation $M_f$ is a free module over $A_f$.
\end{lemma}


\begin{proof} Over the quotient field $K$ of $A$, the $K$-algebra
$B_K = K\otimes _AB$ is of finite type, and $M_K = K\otimes_AM$ is a
finite module over $B_K$. Let $n$ be the dimension of the support
of $M_K$ over $\spec B_K$. We argue by induction on $n$, starting
with $n=-1$ which is the case when $M_K =0$. In this case, as
$K\otimes _AM = S^{-1}M$ where $S = A - \{ 0\}$, each $v\in M$
is annihilated by some non-zero element of $A$. Taking a finite
generating set, and a common multiple of corresponding annihilating
elements, we see there exists an $f\ne 0$ in $A$ with $fM=0$.
Hence $M_f =0$, proving the lemma when $n=-1$.

Now let $n\ge 0$, and let the lemma be proved for smaller values.
As $B$ is noetherian and $M$ is assumed to be a finite $B$-module,
there exists a finite filtration
$$0=M_0 \subset \ldots  \subset M_r =M$$
where each $M_i$ is a $B$-submodule of $M$ such that for each $i\ge 1$
the quotient module $M_i/M_{i-1}$ is isomorphic to $B/\pp_i$ for
some prime ideal $\pp_i$ in $B$.

Note that if $0\to M'\to M\to M'' \to 0$ is a short exact
sequence of $B$-modules, and if $f'$ and $f''$ are non-zero elements of
$A$ such that $M'_{f'}$ and $M''_{f''}$ are free over
respectively $A_{f'}$ and $A_{f''}$, then $M_f$ is a free module over $A_f$
where $f=f'f''$. We will use this fact repeatedly.
Therefore it is enough to prove the result when
$M$ is of the form $B/\pp$ for a prime ideal $\pp$ in $B$.
This reduces us to the case where $B$ is a domain and $M=B$.

As by assumption $K\otimes_AB$ has dimension $n\ge 0$ (that is,
$K\otimes_AB$ is non-zero), the map $A\to B$ must be injective.
By Noether normalisation lemma, there exist elements
$b_1,\ldots,b_n \in B$,
such that $K\otimes_AB$ is finite over its subalgebra
$K[b_1,\ldots, b_n]$ and
the elements $b_1,\ldots, b_n$ are algebraically independent
over $K$. (For simplicity of notation, we write $1\otimes b$ simply as $b$.)
If $g\ne 0$ in $A$ is chosen to be a `common denominator' for
coefficients of equations of integral dependence
satisfied by a finite set of
algebra generators for $K\otimes_AB$ over $K[b_1,\ldots, b_n]$,
we see that $B_g$ is finite over $A_g[b_1,\ldots, b_n]$.

Let $m$ be the generic rank of the finite module $B_g$ over the
domain $A_g[b_1,\ldots, b_n]$. Then we have a short exact
sequence of $A_g[b_1,\ldots, b_n]$-modules of the form
$$0\to A_g[b_1,\ldots, b_n]^{\oplus m} \to B_g \to T \to 0$$
where $T$ is a finite torsion module over $A_g[b_1,\ldots, b_n]$.
Therefore, the dimension of the support of $K\otimes_{A_g}T$
as a $K\otimes_{A_g}(B_g)$-module is strictly less than $n$. Hence by
induction on $n$ (applied to the data $A_g$, $B_g$, $T$),
there exists some $h\ne 0$ in $A$ with $T_h$ free over $A_{gh}$.
Taking $f=gh$, the lemma follows from the
above short exact sequence.
\hfill$\square$\end{proof}



The above theorem has the following consequence, which follows
by restricting attention to a non-empty affine open subscheme of $S$.



\begin{theorem}\label{nhq-generic-flatness}\dcref{4.2}\source{nitsure-hilbert-quot:page-0019}\uses{nhq-algebraic-generic-flatness}
Let $S$ be a noetherian and integral scheme.
Let $p: X\to S$ be a finite type morphism, and let $\F$ be a
coherent sheaf of ${\mathcal O}_X$-modules. Then there exists a non-empty
open subscheme $U\subset S$ such that the restriction of $\F$
to $X_U = p^{-1}(U)$ is flat over ${\mathcal O}_U$.
\end{theorem}

\begin{proof}
This follows by restricting attention to a non-empty affine open subscheme of $S$ and applying Lemma \ref{nhq-algebraic-generic-flatness}.
\end{proof}








%\bigskip





\centerline{\large\bf Existence of Flattening Stratification}

\medskip

\begin{theorem}\label{nhq-flattening-stratification}\dcref{4.3}\source{nitsure-hilbert-quot:page-0019}\uses{nhq-generic-flatness,nhq-graded-module-flat,nhq-base-change-without-flatness,nhq-base-change-flat-sheaves}
Let $S$ be a noetherian scheme, and let
$\F$ be a coherent sheaf on the projective space $\PP^n_S$
over $S$. Then the set $I$ of Hilbert polynomials of
restrictions of $\F$ to fibers of $\PP^n_S\to S$
is a finite set. Moreover, for each $f\in I$ there exist
a locally closed subscheme $S_f$ of $S$,
such that the following conditions are satisfied.

\textbf{(i) Point-set: } The underlying set $|S_f|$ of $S_f$ consists of
all points $s\in S$ where the Hilbert polynomial of the
restriction of $\F$ to $\PP^n_s$ is $f$. In particular,
the subsets $|S_f|\subset |S|$ are disjoint,
and their set-theoretic union is $|S|$.

\textbf{(ii) Universal property: }
Let $S'=\coprod S_f$ be the coproduct of the $S_f$, and let
$i:S'\to S$ be the morphism induced by the inclusions
$S_f\hra S$. Then the sheaf $i^*(\F)$ on $\PP^n_{S'}$
is flat over $S'$. Moreover, $i:S'\to S$ has the universal property that
for any morphism $\varphi :T\to S$ the pullback $\varphi^*(\F)$ on
$\PP^n_T$ is flat over $T$ if and only if $\varphi$ factors through
$i:S'\to S$.  The subscheme $S_f$ is
uniquely determined by the polynomial $f$.

\textbf{(iii) Closure of strata: }
Let the set $I$ of Hilbert polynomials be given a total
ordering, defined by putting $f<g$ whenever $f(n) < g(n)$ for all
$n \gg  0$. Then the closure in $S$ of the subset $|S_f|$
is contained in the union of all $|S_g|$ where $f\le g$.
\end{theorem}


\begin{proof} It is enough to prove the theorem for open subschemes of
$S$ which cover $S$, as
the resulting strata will then glue together by their universal property.


\textbf{Special case: } Let $n=0$, so that $\PP^n_S=S$.
For any $s\in S$, the \textbf{fiber} $\F|_s$ of $\F$ over $s$
will mean the pull-back of $\F$ to the subscheme
$\spec \kappa(s)$, where $\kappa(s)$ is the residue field at $s$.
(This is obtained by tensoring the stalk of $\F$ at $s$ with the
residue field at $s$, both regarded as ${\mathcal O}_{S,s}$-modules.)
The Hilbert polynomial of the
restriction of $\F$ to the fiber over $s$ is the degree $0$
polynomial $e \in \Q[\lambda]$, where $e = \dim_{\kappa(s)}\F|_s$.

By Nakayama lemma, any basis of $\F|_s$ prolongs to a neighbourhood $U$ of $s$
to give a set of generators for $F|_U$. Repeating this argument,
we see that there exists a smaller neighbourhood $V$ of $s$ in which
there is a right-exact sequence
$${\mathcal O}_V^{\oplus m} \stackrel{\psi}{\to} {\mathcal O}_V^{\oplus e}
\stackrel{\phi}{\to} \F \to 0$$
Let $I_e\subset {\mathcal O}_V$ be the ideal sheaf
formed by the entries of the $e\times m$ matrix
$(\psi_{i,j})$ of the homomorphism
${\mathcal O}_V^{\oplus m} \stackrel{\psi}{\to} {\mathcal O}_V^{\oplus e}$.
Let $V_e$ be the closed subscheme of $V$ defined by $I_e$.
For any morphism of schemes
$f: T\to V$, the pull-back sequence
$${\mathcal O}_T^{\oplus m} \stackrel{f^*\psi}{\to} {\mathcal O}_T^{\oplus e}
\stackrel{f^*\phi}{\to}f^* \F \to 0$$
is exact, by right-exactness of tensor products.
Hence the pull-back $f^*\F$ is a locally free ${\mathcal O}_T$-module of
rank $e$ if and only if $f^*\psi =0$, that is,
$f$ factors via the subscheme $V_e\hra V$ defined by the vanishing of all
entries $\psi_{i,j}$.
Thus we have proved assertions \textbf{(i)} and \textbf{(ii)} of the theorem.


As the rank of the matrix $(\psi_{i,j})$ is lower semi-continuous,
it follows that the function $e$ is upper semi-continuous, which
proves the assertion \textbf{(iii)} of the theorem, completing its proof
when $n=0$.


\textbf{General case: } We now allow the integer $n$ to be arbitrary.
The idea of the proof is as follows:
We show the existence of a stratification
of $S$ which is a `g.c.d.' of the flattening
stratifications for direct images $\pi_*\F(i)$ for all $i \ge N$ for some
integer $N$ (where the flattening
stratifications for $\pi_*\F(i)$ exist by case $n=0$ which we have
treated above). This is the desired flattening stratification of $\F$
over $S$, as follows from Lemma \ref{nhq-graded-module-flat}.


As $S$ is noetherian, it is a finite union of irreducible components,
and these are closed in $S$.
Let $Y$ be an irreducible component of $S$, and let $U$ be the
non-empty open subset of $Y$ which consists of all points which do not
lie on any other irreducible component of $S$.
Let $U$ be given the reduced subscheme
structure. Note that this makes $U$ an integral scheme, which is a
locally closed subscheme of $S$.
By Theorem \ref{nhq-generic-flatness} on generic flatness,
$U$ has a non-empty open subscheme $V$ such that
the restriction of $\F$ to $\PP^n_V$ is flat over ${\mathcal O}_V$.
Now repeating the argument with $S$ replaced by its reduced
closed subscheme $S-V$,
it follows by noetherian induction on $S$ that there exist
finitely many reduced, locally closed, mutually disjoint subschemes
$V_i$ of $S$ such that set-theoretically
$|S|$ is the union of the $|V_i|$ and
the restriction of $\F$ to $\PP^n_{V_i}$ is flat over ${\mathcal O}_{V_i}$.
As each $V_i$ is a noetherian scheme, and as the Hilbert polynomials
are locally constant for a flat family of sheaves, it follows
that only finitely many polynomials occur in $V_i$
in the family of Hilbert polynomials $P_s(m) = \chi(\PP^n_s,\F_s(m))$
as $s$ varies over points of $V_i$. This allows us to conclude
the following:

\medskip

\begin{claim}\label{nhq-finitely-many-hilbert-polynomials}\source{nitsure-hilbert-quot:page-0020}\uses{nhq-generic-flatness}
\textbf{(A) } Only finitely many distinct Hilbert polynomials
$P_s(m) = \chi(\PP^n_s,\F_s(m))$ occur, as $s$
varies over all of $S$.

\end{claim}
\medskip

By the semi-continuity theorem applied to the flat families
$\F_{V_i} = \F|_{\PP^n_{V_i}}$ para\-met\-rised by the
finitely many noetherian schemes
$V_i$, we get the following:



\medskip
\begin{claim}\label{nhq-uniform-vanishing-higher-cohomology}\source{nitsure-hilbert-quot:page-0020}\uses{nhq-base-change-flat-sheaves}
\textbf{(B) } There exists an integer $N_1$ such that
$R^r\pi_*\F(m)=0$ for all $r\ge 1$ and $m\ge N_1$,
and moreover $H^r(\PP^n_s,\F_s(m)) = 0$
for all $s\in S$.


\end{claim}
\medskip



For each $V_i$, by Lemma \ref{nhq-base-change-without-flatness}
there exists an
integer $r_i \ge N_1$ with the property
that for any $m\ge r_i$ the base change homomorphism
$$(\pi_*\F(m))|_{V_i} \to {\pi_i}_*\F_{V_i}(m)$$
is an isomorphism, where $\F_{V_i}$ denotes the restriction of
$\F$ to $\PP^n_{V_i}$, and $\pi_i : \PP^n_{V_i}\to V_i$ the projection.
As the higher cohomologies of all fibers (in particular, the first
cohomology) vanish by \textbf{(B)},
it follows by semi-continuity theory for the flat family $\F_{V_i}$
over $V_i$ that for any $s\in V_i$
the base change homomorphism
$$({\pi_i}_*\F_{V_i}(m))|_s \to  H^0(\PP^n_s,\F_s(m))$$
is an isomorphism for $m\ge r_i$. Taking $N$ to be the maximum of
all $r_i$ over the finitely many non-empty $V_i$, and composing the
above two base change isomorphisms, we get the following.


\medskip

\begin{claim}\label{nhq-base-change-direct-image}\source{nitsure-hilbert-quot:page-0021}\uses{nhq-base-change-without-flatness,nhq-uniform-vanishing-higher-cohomology}
\textbf{(C) } There exists an integer $N \ge N_1$ such that
the base change homomorphism
$$(\pi_*\F(m))|_s \to  H^0(\PP^n_s,\F_s(m))$$
is an isomorphism for all $m\ge N$ and $s\in S$.



\end{claim}
\medskip


\textbf{Note } We now forget the subschemes $V_i$ but retain the facts
\textbf{(A)},
\textbf{(B)},
\textbf{(C)}
which were proved using the $V_i$.

\medskip


Let $\pi : \PP^n_S \to S$ denote the projection. Consider the
coherent sheaves $E_0, \ldots, E_n$ on $S$, defined by
$$E_i = \pi_*\F(N+i)\mbox{ for } i=0,\ldots,n.$$
By applying the special case of of the theorem (where the relative dimension
$n$ of $\PP^n_S$ is $0$) to the sheaf $E_0$ on $\PP^0_S=S$, we get
a stratification $(W_{e_0})$ of $S$
indexed by integers $e_0$, such that for any morphism
$f: T\to S$ the pull-back $f^*E_0$ is a locally free ${\mathcal O}_T$-module
of rank $e_0$ if and only if $f$ factors via $W_{e_0} \hra S$.
Next, for each stratum $W_{e_0}$, we take the
flattening stratification $(W_{e_0,e_1})$ for $E_1|_{W_{e_0}}$, and so on.
Thus in $n+1$ steps, we obtain finitely many
locally closed subschemes
$$W_{e_0,\ldots,e_n} \subset S$$
such that for any morphism $f:T\to S$ the pull-back
$f^*E_i$ for $i=0,\ldots,n$ is a locally free ${\mathcal O}_T$-modules of
of constant rank $e_i$ if and only if $f$ factors
via $W_{e_0,\ldots,e_n} \hra S$.



For any integer $N$ and $n$ where $n\ge 0$,
there is a bijection from the set of numerical polynomials
$f\in \Q[\lambda]$ of degree $\le n$ to the set $\Z^{n+1}$,
given by
$$f \mapsto (e_0,\ldots,e_n)\mbox{ where }e_i = f(N+i).$$
Thus, each tuple $(e_0,\ldots,e_n) \in \Z^{n+1}$ can be uniquely
replaced by a numerical polynomial
$f\in \Q[\lambda]$ of degree $\le n$, allowing us to re-designate
$W_{e_0,\ldots,e_n} \subset S$ as $W_f \subset S$.



Note that at any point $s\in S$,
by \textbf{(B)} we have
$H^r(\PP^n_s,\F_s(m)) = 0$ for all $r\ge 1$ and $m\ge N$.
The polynomial $P_s(m) = \chi(\PP^n_s,\F_s(m))$ has degree $\le n$,
so it is determined by its $n+1$ values $P_s(N),\ldots, P_s(N+n)$.
This shows that at any point $s\in W_f$, the Hilbert polynomial
$P_s(m)$ equals $f$. The desired locally closed subscheme $S_f \subset S$,
whose existence is asserted by the theorem,
will turn out to be a certain closed subscheme $S_f \subset W_f$ whose
underlying subset is all of $|W_f|$. The scheme structure of
$S_f$ (which may in general differ from that of $W_f$) is defined as follows.

For any $i\ge 0$ and $s\in S$, the base change homomorphism
$$(\pi_*\F(N+i))|_s \to  H^0(\PP^n_s,\F_s(N+i))$$
is an isomorphism by statement \textbf{(C)}.
Hence each $\pi_*\F(N+i)$ has fibers of constant rank $f(N+i)$
on the subscheme $W_f$. However, this does not mean $\pi_*\F(N+i)$
restricts to a locally constant sheaf of rank $f(N+i)$. But it means that
$W_f$ has a closed subscheme $W_f^i$,
whose underlying set is all of $|W_f|$, such that
$\pi_*\F(N+i)$ is locally free of rank $f(N+i)$ when restricted to
$W_f^{(i)}$, and moreover has the property that any base-change
$T\to S$ under which $\pi_*\F(N+i)$ pulls back to a
locally free sheaf of rank $f(N+i)$ factors via $W_f^i$.
The scheme structure of $W_f^{(i)}$ is defined by a coherent ideal sheaf
$I_i \subset {\mathcal O}_{W_f}$.
Let $I\subset {\mathcal O}_{W_f}$ be the sum of the $I_i$ over $i\ge 0$.
By noetherian condition, the increasing sequence
$$I_0 \subset I_0 + I_1 \subset I_0 + I_1 + I_2 \subset \ldots $$
terminates in finitely many steps, showing $I$ is again a coherent
ideal sheaf. Let $S_f\subset W_f$ be the closed subscheme defined by the
ideal sheaf $I$. Note therefore that $|S_f| = |W_f|$
and for all $i\ge 0$, the sheaf $\pi_*\F(N+i)$ is locally free of
rank $f(N+i)$ when restricted to $S_f$.



It follows that from their definition that
the $S_f$ satisfy property (i) of the theorem.

We now show that the morphism $\coprod_f S_f \,\to \,S$
indeed has the property (ii) of the theorem.
By Lemma \ref{nhq-base-change-without-flatness},
there exists some $N' \ge N$ such that for all $i\ge N'$,
the base-change
$(\pi_*\F(i))|_{S_f} \to (\pi_{S_f})_*\F_{S_f}(i)$ is an isomorphism for
each $S_f$. Therefore
$\F_{S_f}$ is flat over $S_f$ by Lemma \ref{nhq-graded-module-flat},
as the direct images $\pi_*\F(i)$ for all $i \ge N'$ are locally
free over $S_f$. Conversely, if $\phi : T\to S$ is a morphism such
that $\F_T$ is flat, then the Hilbert polynomial is locally
constant over $T$. Let $T_f$ be the open and closed subscheme of $T$
where the Hilbert polynomial is $f$. Clearly, the set
map $|T_f| \to |S|$ factors via $|S_f|$. But as the direct images
${\pi_T}_*\F_T(i)$ are locally free of rank $f(i)$ on $T_f$, it
follows in fact that the schematic morphism $T_f \to S$ factors
via $S_f$, proving the property (ii) of the theorem.

As by \textbf{(A)} only finitely many polynomials $f$ occur, there
exists some $p \ge N$ such that for any two polynomials $f$ and $g$
that occur, we have $f < g$ if and only if $f(p) < g(p)$.
As $S_f$ is the flattening stratification for $\pi_*\F(p)$,
the property (iii) of the theorem follows from the corresponding property in
the case $n=0$, applied to the sheaf $\pi_*\F(p)$ on $S$.

This completes the proof of the theorem.  \hfill$\square$\end{proof}


%\bigskip
